\section{怎样写比较简单的记叙文}

记叙文是主要用来写人记事的一种文体。就初学写作的初中学生来说，应该先从比较简单的记叙文练起。说它“比较简单”，是指写作的内容只限于一人、一事、一物、一景；写作的要求只限于把要写的人、事、物、景，具体、完整、清楚、明白地表达出来，显示出一个明确的中心意思。

在动手作文之前，必须先做好审题、构思和编写提纲等准备工作。俗话说：磨刀不误砍柴工。只有做好了这些准备工作，大致确定了写什么和怎么写之后，才能胸有成竹，行文流畅。关于这些准备工作的具体做法，我们在前面已作了说明。现在只谈谈怎样具体来写比较简单的写人记事的记叙文。

\textbf{1. 学会写好一个人}

老师常常出这样一类作文题：《记三好学生×××》、《我最尊敬的一个人》、《我的好友×××》、《我最难忘的一位老师》、《亲爱的妈妈》、《×× --- 我们学习的榜样》、《我》、《××速写》等。这一类题目，都是叫我们写一个比较了解、熟悉的人物，通过对这个人物的叙述和描写来反映现实生活的一个方面，表达文章的中心思想。

我们怎样写好这个人呢？首先得想想这个人的思想、感情和个性怎么样，要找出这个人的特点。再想想这个人的特点通过哪一个比较完整的事件能够突出地表现出来，或者通过哪一、两个片断能够突出地表现出来。例如写《记三好学生×××》，假如已确定他的特点是不仅严格要求自己，而且勇于帮助别人，那么他有哪些事实可以表现这个特点呢？也许你会想到他在数学考了98分时，怎样严格订正错误或者怎样进一步思索更好的解题方法；也许你会想到他在大家忙于复习功课、无暇打扫教室时，他却悄悄地一个人做了；也许你会想到某个同学生病而耽误了功课时，他却风雨无阻地为那位同学补习功课，等等。如果你认为他给同学补习功课的事有头有尾，可以构成一个生活小故事，而且这件事最能表现他的思想品质和性格特点，那就集中写这件事，写出事情的发生、发展和结束的全过程，去刻画这位三好同学的鲜明性格，表现歌颂他的思想境界的主题。如果你认为只有把这几个生活片断组合起来，从几个方面写才能突出表现这位三好同学的思想境界，当然也是可以的，只是这几个生活片断不必写成一个个有头有尾的故事，能写得相对完整又比较具体就行。需要注意的是，所写的事一定要为表现人物的思想品质和性格特点服务，而且要写得具体生动，切忌笼统的叙述和过多的概括介绍，尤其不宜以空论代替具体的叙述。

为了使人物的性格特征、内心世界活灵活现地展示出来，除了具体生动地叙述事实，还需要对人物的外貌、语言、动作、心理以及所处的环境进行恰当的描写。描写外貌，要善于用简练的语言勾勒出人物的特点，正象鲁迅先生所说：“要极省俭地画出一个人的特点，最好是画他的眼睛。”一定要以自己的观察为依据，使外貌描写切合人物的年龄、身分和性格，不要想当然地生搬硬套习见的词语。语言描写，既要合乎人物的身分、经历、个性，又要口语化，给人以真实感。行动描写要具体、细致、生动，有鲜明的个性，其中更要注意选择准确、形象的动词，以刻画人物的个性和精神面貌。心理描写要合乎特定环境下的心理状态，具有个性特点，文字不可冗长，使人读了感到沉闷。环境描写则要符合人物的心理状态，能有力地渲染环境气氛，烘托人物性格，切切不可为写景而写景。

我们只要能抓住一个人的思想性格的特点，找到并且善于叙述最能表现这种特点的事实，再借助于各种描写人物的手法，就能把一个人的形象写好，从而深刻地表现出文章的中心意思。

\textbf{2. 学会写好一件事。}

老师还常常出这样一类作文题：《一件有意义的小事》、《难忘的一堂课》、《拔河比赛》、《记一次参观活动》、《瞻仰烈士陵园》等。这类题目，或是叫我们写出生活中的一个小故事，或是叫我们记叙一场体育比赛、一堂课、一次参观访问、一次社会活动，等等。写这些小故事或者活动，当然都离不开写人，但是它们不像上面所说的那样主要去写一个人，而是主要去写一件事或一次活动的过程。我们要通过对一件事或一次活动的叙述和描写，来反映社会生活，表达文章的中心思想，而不一定对人物的思想性格进行具体的刻画。由于这样的一件事或一次活动，展开的生活面一般比较窄，事件延续的时间比较短，发展变化比较简单，涉及的人也比较少（即使写群众性的活动，也只能点一下比较突出的一两个人），所以仍然属于比较简单的记叙文。

我们怎样写好一件事（或一次活动）呢？

首先，要把一件事的发生、发展和结束的整个过程想得清楚，记得完整，并且按照一定的时间顺序和空间位置把事件交代得明明白白，有条不紊。初学写作的人往往会犯交代不清的毛病：有时因为考虑不周或行文草率，使情节当中出现漏洞；有时对事情的前因后果处理得不合情理或缺少必要的说明，有时忘了点明时间的推移或地点的转换，等等。这些毛病出现任何一个，都会使读者产生疑问，削弱了文章的真实性。

其次，要充分描述事件的具体情况，让事实本身说话，去显示中心思想，而极力避免抛开具体事实说空话。记事要记得具体，是记叙文写作的特点。为了做到这一点，就要求我们平时注意观察生活实际，善于抓住事件的要点和细节；在作文时，应该选用那些亲身经历、感受较深的事；并在写作过程中，多问自己几个“怎么样”，把要写的事件的有关情况想全面，想细致，想得能达到历历如在眼前的境地，写出来的文章就一定会具体动人。初学写作的人，往往既不能把一件事从头到尾详细展开，又把事件进程中许多应该展开的情节匆匆几笔带过，使事件本身只剩下个粗率的轮廓。更有甚者，用大段大段的议论或抒情去代替事件本身的记叙和描写，使一点点具体的材料淹没在空话的汪洋大海里。这样，或者使文章内容空洞抽象，或者使文章简直不象篇记叙文了，我们一定要避免。

再次，要紧紧围绕中心，精心选取那些足以表达中心思想的材料，把它集中地突出地写出来，而与中心思想无关的材料一概放到一边去。对那些最能表现中心思想的材料，要写得具体而详尽，对那些只起陪衬和连贯作用的材料，就应写得概括而简略。初学写作的人，往往会抓不住中心，想到哪里就写到哪里，使读者看不出写的中心意思究竟是什么，有的是既想写这个，又想写那个，这个写一点，那个写一点，头绪多，内容杂，看不出哪个是在表达中心意思，还有的虽然有了要表达的中心意思，却不围绕中心去选材和组织材料，从而把无关紧要的话说得太多，有关中心思想的话只寥寥几笔，结果喧宾夺主，没有把中心思想很好地表达出来。为了做到记事中心突出，详略得当，必须在编写提纲时就明确所要表达的中心思想，并按照中心思想选材和组织材料，确定详略安排，然后依照提纲写作，就可避免出现没有中心、详略不当的毛病了。

最后，要讲究表达方法，力求构思奇巧，使故事的情节发展曲折有致，引人入胜。这是写好一件事的较高要求。初学写作的人，在明确了中心思想和有了一个能表现中心思想的故事之后，往往是按着事件的先后顺序从头一直写到结束，象记一篇流水帐，让人一览无余，显得平淡无奇，读来乏味。为了使文章能紧紧抓住读者的心，就得在叙事方法的灵活多变方面多下功夫：或者制造悬念，引人急于知道下情；或者采取倒叙，让人很想了解根由；或者使情节发展环环相扣，层层推进，使故事波澜起伏，引人入胜；或者巧取叙事角度，用侧面描写来映衬正面描写，用小事细物来写出重大情节，等等。一句话，就是让读者感到出乎意料之外，又在情理之中。这样来记一件事，怎能会不生动、不吸引人看呢？

\textbf{3. 学会安排好比较简单的记叙文的结构。}

记叙文的结构能够反映出作者对于情节发展变化全过程的构思。结构安排得好，能使作者的思想观点和写作材料找到一种最恰当又最有力的表现形式。有人说，如果把主题看成是文章的“灵魂”，材料是文章的“血肉”，那么，结构就是文章的“骨骼”。这个比喻是很有道理的。试想，没有坚实匀称的骨骼，血肉和灵魂就不能有所依附，怎能成为一个活生生的能够站立行走的真人呢！

要想安排好结构，作者的思路必须清晰、严密才行。因为只有思路清晰、严密，才能合理地划分段落，处理好上下文的过渡和照应，使文章层次分明，条理清楚，前后连贯，首尾一致。

一般说来，文章的一层意思可以相对独立的，就划分为一个段落；如果一层意思又包含几方面的意思，就可以把这几方面的意思再分别划段，以免有的段落非常长，使读者感到眉目不清。有的同学习惯于把一篇作文划分为开头、结尾和中间主干部分三段，很明显，那中间主干部分作为一段来说，肯定内容多，显得臃肿，表意不清，应该适当地再划分一下。写文章的时候，应该力求把每段的意思说得完整、清楚。不要把不该在这一段里说的意思硬往里面塞，也不要一段意思没有说清楚，又扯到另一段去。

就全文来看，段落只是整篇文章的一个组成部分，因此必须使前后段落保持连贯性。为此，就要安排好上文和下文的过渡和照应。过渡是连接层次、段落的桥梁，使各个层次和段落连接成为一个有机的整体。大致说来有以下几种情况需要过渡：从这个情节发展到另一个情节的叙述，如课文《人民的勤务员》，由叙述转入议论和由议论转入叙述，如课文《一件珍贵的衬衫》，由叙述现在的事件插入叙述过去的事件，如课文《红军鞋》。过渡的方法大致有以下几种：

（1）可用设问句过渡，如课文《梁生宝买稻种》用“他为什么不进旅馆呢？”过渡到对梁生宝不住旅馆原因的议论；又用“到哪里过一夜呢？”过渡到梁生宝的回忆。

（2）用关联词语过渡，如课文《第比利斯的地下印刷所》第二段第一层总叙小院子的外景和内部结构以后，用“可是秘密就在这口井里”，过渡到第二层，转入揭示这口井的秘密。

（3）用陈述句或小节过渡，如课文《一件小事》用独立的一个小节“但有一件小事，却于我有意义，将我从坏脾气里拖开，使我至今忘记不得。”自然地使文章从议论转入对一件小事的回忆。有的同学写作时常常忘记过渡，一直写下去，致使文章脉络不清，结构松散，应该努力避免。

为了能使段落与段落之间前后映衬，互相呼应，使文章文理顺畅，首尾贯通，还要在作文时注意上下文的照应。照应的方法常见的有两种：一是行文中的照应，二是首尾照应。如课文《驿路梨花》，开头先写在高山密林里看到梨花，又写瑶族老人在困境中看到了梨花，接写原来是一个名叫梨花的小姑娘在为行人做好事，再写自己梦见这位小姑娘在梨花丛中歌唱，后面又写梨花的妹妹象梨花一样在为行人做好事，最后以“驿路梨花处处开”结尾，这样的写法就是行文中的照应，它使文章结构非常严密，一线贯通，有力地表达出中心思想。再如课文《老山界》开头一句写明“我们决定要爬一座三十里高的瑶山——老山界”，最后用“老山界是我们长征中所过的第一座难走的山······”结尾，使文章首尾遥相呼应，成为一个统一的整体。有不少文章在行文中常与题目照应，这种照应题目的地方往往起点题的作用。同学们在作文时，有时前面写的活，后面忘了顾及，使前面的话落了空；有时前面没有作交代，后文却突然冒出一件事物来，令人茫然不解。因此，我们在写文章时，开头的描写一定要为情节的发展提供线索，后面要说的话，前面先要交代一下，即所谓打下“伏笔”。前面有因，后面一定要有果；前面有悬念，后面一定要有结局。这样才能前后相承，首尾照应，脉络细密，引人入胜。

总起来说，写人要写出人的形象特点和精神面貌，记事要具体地叙述和描写事件本身的发展和变化，并且要灵活运用各种写人记事的方法，力求能让读者如见其人，如临其境，情节起伏，结构分明。在必要的画龙点睛的地方可以进行简要的议论或抒情，但是决不可过长过滥地发表议论或抒发感情。如果议论、抒情过多，记叙过少，那就不能叫做记叙文，至少是不大象记叙文了。
\newpage